\subsection{\texorpdfstring{Properties of the ring \(A^{(d)}\)}{Properties of the ring A(d)}}

Let \(\mathbf{A} = \bigoplus_{i\in \mathbf{Z}}\mathbf{A}_i\) be a graded ring and \(\mathbf{M} = \bigoplus_{i\in \mathbf{Z}}\mathbf{M}_i\) a graded A-module; for every ordered pair of integers \((d, k)\) such that \(d \geqslant 1\) and \(0 \leqslant k \leqslant d - 1\) , set

\[
\mathbf {A} ^ {(d)} = \bigoplus_ {i \in \mathbf {Z}} \mathbf {A} _ {i d}, \quad \mathbf {M} ^ {(d, k)} = \bigoplus_ {i \in \mathbf {Z}} \mathbf {M} _ {i d + k}.
\]

Clearly \(\mathbf{A}^{(d)}\) is a graded subring of A and \(\mathbf{M}^{(d,k)}\) a graded \(\mathbf{A}^{(d)}\) -module; moreover, if N is a graded submodule of M, \(\mathbf{N}^{(d,k)}\) is a graded sub- \(\mathbf{A}^{(d)}\) -module of \(\mathbf{M}^{(d,k)}\) . We shall write \(\mathbf{M}^{(d)}\) instead of \(\mathbf{M}^{(d,0)}\) ; for each \(d \geqslant 1\) , M is the direct sum of the \(\mathbf{A}^{(d)}\) -modules \(\mathbf{M}^{(d,k)}\) ( \(0 \leqslant k \leqslant d - 1\) ).

PROPOSITION 2. Let \(\mathbf{A} = \bigoplus_{i\in \mathbf{Z}}\mathbf{A}_i\) be a graded commutative ring with positive degrees and \(\mathbf{M} = \bigoplus_{i\in \mathbf{Z}}\mathbf{M}_i\) a graded A-module. Suppose that \(\mathbf{A}\) is a finitely generated \(\mathbf{A}_0\) -algebra and \(\mathbf{M}\) a finitely generated A-module. Then, for every ordered pair \((d,k)\) of integers such that \(d\geqslant 1,0\leqslant k\leqslant d - 1\) :

\begin{enumerate}
\def\labelenumi{(\roman{enumi})}
\item
  \(\mathbf{A}^{(d)}\) is a finitely generated \(\mathbf{A}_0\) -module.
\item
  \(\mathbf{M}^{(d,k)}\) is a finitely generated \(\mathbf{A}^{(d)}\) -module.
\end{enumerate}

Let us show that A is a finitely generated \(\mathbf{A}^{(d)}\) -module. Let \((a_{i})_{1\leqslant i\leqslant s}\) be a system of generators of the \(\mathbf{A}_0\) -algebra A consisting of homogeneous elements. The elements of A (finite in number) of the form \(a_1^{\alpha_1}a_2^{\alpha_2}\ldots a_s^{\alpha_s}\) such that \(0\leqslant \alpha_{i}\leqslant d\) for \(1\leqslant i\leqslant s\) constitute a system of generators of the \(\mathbf{A}^{(d)}\) -module A; for every system of integers \(n_i\geqslant 0\) ( \(1\leqslant i\leqslant s\) ), there are positive integers \(q_{i}, r_{i}\) such that \(n_i = q_id + r_i\) where \(r_i < d\) ( \(1\leqslant i\leqslant s\) ); then

\[
a _ {1} ^ {n _ {1}} a _ {2} ^ {n _ {2}} \dots a _ {s} ^ {n _ {s}} = (a _ {1} ^ {q _ {1}} \dots a _ {s} ^ {q _ {s}}) ^ {d} (a _ {1} ^ {r _ {1}} \dots a _ {s} ^ {r _ {s}})
\]

which proves our assertion, for every homogeneous element \(x \in A\) satisfies \(x^{d} \in \mathrm{A}^{(d)}\) . Then, if M is a finitely generated A-module, it is also a finitely generated \(\mathrm{A}^{(d)}\) -module; as M is the direct sum of the \(\mathrm{M}^{(d,k)}\) ( \(0 \leqslant k \leqslant d - 1\) ), each of the \(\mathrm{M}^{(d,k)}\) is a finitely generated \(\mathrm{A}^{(d)}\) -module, which proves (ii).

Let us apply the above to the graded A-module \(m = \bigoplus_{i \geqslant 1} A_i\) , which is finitely generated by the Corollary to Proposition 1 of no. 2; it is seen that \(m^{(d)}\) is a finitely generated \(A^{(d)}\) -module; therefore (no. 2, Corollary to Proposition 1) \(A^{(d)}\) is a finitely generated \(A_0\) -algebra.

LEMMA 1. Let A be a graded commutative ring such that \(\mathbf{A} = \mathbf{A}_0[\mathbf{A}_1]\) , M a graded A-module and \((y_{\lambda})_{\lambda \in \mathbf{L}}\) a system of homogeneous generators of M such that \(\deg(y_{\lambda}) \leqslant n_0\) for all \(\lambda \in \mathbf{L}\) . Then, for all \(n \geqslant n_0\) and all \(k \geqslant 0\) , \(\mathbf{M}_{n+k} = \mathbf{A}_k \cdot \mathbf{M}_n\) .

Let \(n \geqslant n_0\) and \(x \in \mathbf{M}_{n+1}\) . Since the \(y_\lambda\) generate \(\mathbf{M}\) , there exists a family \((a_\lambda)_{\lambda \in \mathbf{L}}\) of elements of \(\mathbf{A}\) of finite support such that \(x = \sum_{\lambda} a_\lambda y_\lambda\) ; we may further suppose that each \(a_\lambda\) is homogeneous and of degree \(n + 1 - \deg(y_\lambda)\) (replacing it if need be by its homogeneous component of that degree). As

\(A = A_{0}[A_{1}]\) and \(\deg(a_{\lambda}) > 0\) , each \(a_{\lambda}\) is a sum of elements of the form \(bb'\) where \(b \in A_{1}\) , \(b' \in A\) , whence \(x \in A_{1}M_{n}\) . Then \(M_{n+1} = A_{1}M_{n}\) , whence \(M_{n+k} = A_{k}M_{n}\) by induction on k.

LEMMA 2. Let A be a graded commutative ring such that \(\mathbf{A} = \mathbf{A}_0[\mathbf{A}_1]\) and let \(\mathbf{S} = \bigoplus_{i\geqslant 0}\mathbf{S}_i\) be a graded commutative A-algebra with positive degrees, which is a finitely generated A-module. Then there exists an integer \(n_0\geqslant 0\) such that:

\begin{enumerate}
\def\labelenumi{(\roman{enumi})}
\item
  For \(n \geqslant n_0\) and \(k \geqslant 0\) , \(\mathrm{S}_{n+k} = \mathrm{S}_k \cdot \mathrm{S}_n\) .
\item
  For \(d \geqslant n_0\) , \(\mathbf{S}^{(d)} = \mathbf{S}_0[\mathbf{S}_d]\) .
\end{enumerate}

By Lemma 1 there exists an integer \(n_0 \geqslant 0\) such that, for \(n \geqslant n_0\) and \(k \geqslant 0\) , \(S_{n+k} = A_kS_n\) , whence a fortiori \(S_{n+k} = S_kS_n\) , which establishes (i). Then, for \(d \geqslant n_0\) and \(m > 0\) , \(S_{md} = (S_d)^m\) as follows by induction on \(m\) applying (i); this establishes (ii).

PROPOSITION 3. Let \(\mathbf{R} = \bigoplus_{i\geqslant 0}\mathbf{R}_i\) be a graded commutative ring with positive degrees which is a finitely generated \(\mathbf{R}_0\) -algebra. There exists an integer \(e\geqslant 1\) such that \(\mathbf{R}^{(me)} = \mathbf{R}_0[\mathbf{R}_{me}]\) for all \(m\geqslant 1\) .

Let \((x_{j})_{1\leqslant j\leqslant s}\) be a system of homogeneous generators of the R \(_{0}\) -algebra R, whose degrees are \(\geqslant 1\) . Let \(h_{j} = \deg(x_{j})\) , let \(q\) be a common multiple of the \(h_{j}\) and let us write \(q_{j} = q/h_{j}\) for \(1 \leqslant j \leqslant s\) ; the elements \(x_{j}^{q_{j}}\) are then all of degree \(q\) . Let B be the sub-R \(_{0}\) -algebra of R generated by the \(x_{j}^{q_{j}}\) ; it is a graded subalgebra of R and B \(_{i} = 0\) if \(i\) is not a multiple of \(q\) . Let A (resp. S) be the graded ring whose underlying ring is B (resp. R \(^{(q)}\) ) and whose graduation consists of the A \(_{i} =\) B \(_{iq}\) (resp. S \(_{i} =\) R \(_{iq}\) ). Then A = A \(_{0}\) {[}A \(_{1}\) {]} by definition of B. Consider the elements of R (finite in number) of the form \(x_{1}^{\alpha_{1}}x_{2}^{\alpha_{2}}\ldots x_{s}^{\alpha_{s}}\) , where \(0 \leqslant \alpha_{j} \leqslant q_{j}\) and \(\alpha_{1}h_{1} + \cdots + \alpha_{s}h_{s} \equiv 0 \pmod{q}\) ; let us show that they generate the B-module R \(^{(q)}\) . It is sufficient to prove that every element of R \(^{(q)}\) of the form \(x_{1}^{n_{1}}x_{2}^{n_{2}}\ldots x_{s}^{n_{s}}\) is a B-linear combination of the above elements. Now, there exist positive integers \(k_{j}, r_{j}\) such that \(n_{j} = k_{j}q_{j} + r_{j}\) where \(r_{j} < q_{j}\) ( \(1 \leqslant j \leqslant s\) ); then

\[
x _ {1} ^ {n _ {1}} x _ {2} ^ {n _ {2}} \dots x _ {s} ^ {n _ {s}} = (x _ {1} ^ {q _ {1}}) ^ {k _ {1}} \dots (x _ {s} ^ {q _ {s}}) ^ {k _ {s}}. (x _ {1} ^ {r _ {1}} \dots x _ {s} ^ {r _ {s}})
\]

and by hypothesis \(\sum_{j=1}^{s} n_j h_j \equiv 0 \pmod{q}\) , hence \(\sum_{j=1}^{s} r_j h_j \equiv 0 \pmod{q}\) ; as the \(x_j^{q_j}\) belong to B by definition, this proves our assertion. As S is a finitely generated A-module, Lemma 2 can be applied: there exists \(n_0\) such that for \(d \geqslant n_0\) , \(\mathrm{S}^{(d)} = \mathrm{S}_0[\mathrm{S}_d]\) and hence \(\mathrm{R}^{(qd)} = \mathrm{R}_0[\mathrm{R}_{qd}]\) for \(d \geqslant n_0\) . The proposition follows with \(e = q n_0\) .

